% !PRINTED 159
\weberpage{159}

Im Gebiete der binären Substitutionen ist die Substitution $A$:
\[
(y_1, y_2) = \begin{pmatrix} a, & b \\ c, & d \end{pmatrix} (x_1, x_2)
\]
als Substitution der Verhältnisse aufgefasst, gleichbedeutend mit der \gesperrt{linearen gebrochenen Substitution}:
\[
\eta = \frac{a\,\xi + b}{c\,\xi + d},
\]
wenn $\xi = x_1 : x_2$ und $\eta = y_1 : y_2$ gesetzt wird; kommt eine zweite Substitution $A'$:
\[
(z_1, z_2) = \begin{pmatrix} a', & b' \\ c', & d' \end{pmatrix} (y_1, y_2),
\]
oder
\[
\zeta = \frac{a'\eta + b'}{c'\eta + d'}
\]
hinzu, so erhält man die zusammengesetzte Substitution
\[
A'' = A'A,
\]
durch welche $\zeta$ durch $\xi$ ausgedrückt wird:
\[
\zeta = \frac{a''\xi + b''}{c''\xi + d''},
\]
nach den Regeln der Composition in der Form
\[
\begin{pmatrix} a'', & b'' \\ c'', & d'' \end{pmatrix} = \begin{pmatrix} a', & b' \\ c', & d' \end{pmatrix} \begin{pmatrix} a, & b \\ c, & d \end{pmatrix} = \begin{pmatrix} a'a + b'c, & a'b + b'd \\ c'a + d'c, & c'b + d'd \end{pmatrix}.
\]
Ist eine solche Substitution multiplicativ, hat sie also die Form
\[
\begin{pmatrix} a, & 0 \\ 0, & d \end{pmatrix},
\]
so werden wir auch das Verhältniss $a : d$ den \gesperrt{Multiplicator} nennen.

\begin{center}
\S.\,39.

\gesperrt{Permutationen als lineare Substitutionen.}
\end{center}

Die Wichtigkeit der linearen Substitutionen und besonders der daraus gebildeten \gesperrt{endlichen} Gruppen für die Algebra ergiebt sich daraus, dass die Permutationsgruppen von $n$ Elementen als specielle Fälle solcher Substitutionsgruppen aufgefasst werden können.
